\chapter*{Resumo}
\addcontentsline{toc}{chapter}{Resumo}
\markboth{Resumo}{Resumo}

Este relatório apresenta o \rat{}, uma plataforma integrada de análise estática e dinâmica para
deteção de cavalos de Troia de acesso remoto (\emph{RAT}) em executáveis \emph{Windows} (e código
\emph{C\#} com âmbito reduzido), desenvolvida no âmbito do projeto final de Licenciatura em
Engenharia Informática na \ubi{}. A solução combina um servidor (\emph{backend}) em \emph{Python},
com \emph{FastAPI}, e uma cadeia de processamento (\emph{pipeline}) de sete fases (metadados de
executáveis portáteis (\emph{PE}), descompilação com \emph{ILSpy}, heurísticas, desofuscação, regras
\emph{YARA}, pontuação de risco (\emph{risk scoring}), análise nativa com \emph{Ghidra} quando
necessário e geração de relatório), uma interface para navegador (\emph{web}) em \emph{React}, com
assistente \emph{Gemini} opcional, uma aplicação de secretária (\emph{desktop}) \emph{WPF} e um
ambiente de execução isolado (\emph{sandbox}) \emph{Hyper-V} com dois caminhos de orquestração
independentes (agente \emph{HTTP} na máquina virtual (\emph{vm-agent}) e \emph{PowerShell Direct}
com verificação de integridade \emph{SHA-256}). A validação assenta em 220 testes automatizados,
métricas substitutas de deteção (\emph{proxy detection}) em 59 perfis sintéticos etiquetados
estratificados por cinco níveis de severidade (precisão, sensibilidade (\emph{recall}), $F_1$ e
exatidão de nível 1{,}00), um conjunto benigno de oito executáveis \emph{PE} compilados analisados
estaticamente (especificidade 1{,}00, zero falsos positivos), metadados públicos de indicadores de
compromisso (\emph{IOC}) do \emph{MalwareBazaar} (sem armazenamento de software malicioso
(\emph{malware}) na máquina anfitriã (\emph{host})), uma medição formal de desempenho
(\emph{benchmark}) estático (mediana 36{,}1~s, $n=5$), perfis sintéticos de pontuação
(\emph{scoring}) alinhados com \filepath{test\_risk\_scoring.py} e execução benigna através de
\filepath{benign-vm-test.exe} no Caminho~B. O sistema foi concebido para explicabilidade e
isolamento local, em detrimento de soluções comerciais de análise na nuvem (\emph{cloud sandboxes}),
destinando-se a uso académico e laboratorial.

\Needspace{4\baselineskip}
\vspace{0.8em}
\noindent\textbf{Palavras-chave:} análise de software malicioso (\emph{malware}); troianos de acesso
remoto (\emph{RAT}); ambiente de execução isolado (\emph{sandbox}); cibersegurança; engenharia de
software; análise estática e dinâmica.
